\section{Advanced Views}

\begin{frame}{Beyond Form and List}
	\begin{itemize}
		\item Form and list/tree are only the beginning.
		\item Advanced views help users analyze and organize records visually.
		\item In this section:
		\begin{itemize}
			\item Kanban
			\item Pivot
			\item Graph
			\item Calendar
			\item Activity
			\item Gantt
			\item Search view reminders for all of them
		\end{itemize}
	\end{itemize}
\end{frame}

\begin{frame}[fragile]{Declaring Multiple View Modes}
\begin{lstlisting}[style=xml]
<record id="action_student" model="ir.actions.act_window">
  <field name="name">Students</field>
  <field name="res_model">student.student</field>
  <field name="view_mode">kanban,list,form,pivot,graph,calendar,activity</field>
</record>
\end{lstlisting}
	\begin{itemize}
		\item The order in \texttt{view\_mode} controls the switcher order.
		\item Each mode needs a corresponding view architecture.
	\end{itemize}
\end{frame}

\begin{frame}[fragile]{Kanban View}
	\begin{itemize}
		\item Kanban shows records as cards, often grouped in columns.
\begin{lstlisting}[style=xml]
<kanban default_group_by="classroom_id" class="o_kanban_small_column">
  <field name="name"/>
  <field name="classroom_id"/>
  <field name="age"/>
  <templates>
	<t t-name="card">
	  <div class="fw-bold"><field name="name"/></div>
	  <div>Age: <field name="age"/></div>
	</t>
  </templates>
</kanban>
\end{lstlisting}
		\item Great for pipelines, classrooms, stages, and boards.
	\end{itemize}
\end{frame}

\begin{frame}[fragile]{Kanban Tips}
	\begin{itemize}
		\item Use \texttt{default\_group\_by} for column grouping.
		\item Keep cards compact: title + a few key fields.
		\item Include fields used in the template as \texttt{<field/>} declarations.
		\item In Odoo 17, card templates often use \texttt{t-name="card"}.
	\end{itemize}
\end{frame}

\begin{frame}[fragile]{Pivot View}
	\begin{itemize}
		\item Pivot is for interactive spreadsheet-like analysis.
\begin{lstlisting}[style=xml]
<pivot string="Student Analysis">
  <field name="classroom_id" type="row"/>
  <field name="gender" type="col"/>
  <field name="age" type="measure"/>
</pivot>
\end{lstlisting}
		\item \texttt{type="row"} / \texttt{col"} define axes.
		\item \texttt{type="measure"} defines aggregated values.
	\end{itemize}
\end{frame}

\begin{frame}[fragile]{Graph View}
	\begin{itemize}
		\item Graph shows bar/line/pie charts.
\begin{lstlisting}[style=xml]
<graph string="Students by Classroom" type="bar">
  <field name="classroom_id"/>
  <field name="age" type="measure"/>
</graph>
\end{lstlisting}
		\item Useful for dashboards and quick comparisons.
		\item Common types: \texttt{bar}, \texttt{line}, \texttt{pie}.
	\end{itemize}
\end{frame}

\begin{frame}[fragile]{Calendar View}
	\begin{itemize}
		\item Calendar places records on dates.
\begin{lstlisting}[style=xml]
<calendar string="Student Events"
		  date_start="start_datetime"
		  date_stop="end_datetime"
		  color="classroom_id"
		  mode="month">
  <field name="name"/>
  <field name="classroom_id"/>
</calendar>
\end{lstlisting}
		\item Needs at least a start datetime/date field.
	\end{itemize}
\end{frame}

\begin{frame}[fragile]{Activity View}
	\begin{itemize}
		\item Activity view focuses on scheduled activities and follow-ups.
\begin{lstlisting}[style=xml]
<activity string="Student Activities">
  <field name="name"/>
  <field name="activity_state"/>
  <field name="activity_ids"/>
</activity>
\end{lstlisting}
		\item Usually used with \texttt{mail.activity.mixin}.
		\item Excellent for CRM-like follow-up workflows.
	\end{itemize}
\end{frame}

\begin{frame}[fragile]{Gantt View}
	\begin{itemize}
		\item Gantt shows time ranges on a timeline (Enterprise in many Odoo setups).
\begin{lstlisting}[style=xml]
<gantt string="Classroom Schedule"
	   date_start="start_datetime"
	   date_stop="end_datetime"
	   default_group_by="classroom_id"
	   color="teacher_id">
  <field name="name"/>
</gantt>
\end{lstlisting}
		\item Ideal for planning, allocation, and scheduling.
	\end{itemize}
\end{frame}

\begin{frame}[fragile]{Search View Support for Advanced Views}
\begin{lstlisting}[style=xml]
<search>
  <field name="name"/>
  <field name="classroom_id"/>
  <filter string="Adults" name="adults" domain="[('age','&gt;=',18)]"/>
  <group>
	<filter string="Classroom" name="group_classroom"
			context="{'group_by': 'classroom_id'}"/>
	<filter string="Gender" name="group_gender"
			context="{'group_by': 'gender'}"/>
  </group>
</search>
\end{lstlisting}
	\begin{itemize}
		\item Group-by filters are especially important for pivot/graph/kanban.
	\end{itemize}
\end{frame}

\begin{frame}{Choosing the Right Advanced View}
	\begin{itemize}
		\item \textbf{Kanban}: boards / stages / cards
		\item \textbf{Pivot}: analytical tables
		\item \textbf{Graph}: charts and trends
		\item \textbf{Calendar}: date-based events
		\item \textbf{Activity}: follow-up management
		\item \textbf{Gantt}: schedules and allocations
	\end{itemize}
\end{frame}
